\section{Domain Task 1: Phân tích Thị trường và Phân bố không gian}

Câu hỏi nghiệp vụ: ``Khu vực nào tại New York đang có mức giá cho thuê
cao nhất/thấp nhất, và cơ cấu loại phòng ở từng quận phân bố như thế nào?''

Mục đích: Giúp hiểu rõ bức tranh giá cả và cấu trúc nguồn cung theo từng
vùng địa lý để định vị phân khúc thị trường.

\subsection{Biểu đồ 1: Bản đồ phân bố giá trung bình theo khu vực (Choropleth Map)}

\subsubsection*{A. Thiết kế Idiom}

\begin{longtable}{@{} p{2.8cm} p{12cm} @{}}
\toprule
\textbf{Đặc điểm} & \textbf{Chi tiết} \\
\midrule
\endhead

\bottomrule
\endfoot

\textbf{Idiom} & Choropleth Map (Bản đồ phân bố sắc độ) \\
\addlinespace[0.2cm]

\textbf{What} & Tọa độ (Longitude/Latitude): Q, Khu vực (Neighbourhood): C, Giá trung bình (AVG Price): Q, Nhãn ngoại lệ (Price Is Outlier): C \\
\addlinespace[0.2cm]

\textbf{Why} & produce (derive) $\rightarrow$ explore $\rightarrow$ locate $\rightarrow$ summarize, compare
\begin{itemize}[nosep, leftmargin=*, topsep=0.1cm]
    \item Tạo ra (derive) giá trị trung bình của từng khu vực để so sánh.
    \item Khám phá (Explore) bức tranh tổng quan về phân bố giá thuê trên toàn bộ không gian địa lý của 5 quận New York.
    \item Xác định vị trí (Locate) các cụm khu vực đắt đỏ nhất hoặc bình dân nhất.
    \item Tóm tắt (Summarize) và so sánh (Compare) mức chênh lệch giá trị bất động sản/lưu trú giữa vùng trung tâm và vùng lân cận.
\end{itemize} \\
\addlinespace[0.2cm]

\textbf{How} & \textbf{Encode:}
\begin{itemize}[nosep, leftmargin=*, topsep=0.1cm]
    \item Mark: Area (vùng bản đồ)
    \item Channel:
    \begin{itemize}[nosep, leftmargin=0.5cm]
        \item PosX, PosY: Vị trí địa lý (Kinh độ, Vĩ độ) thực tế.
        \item Color (Luminance/Độ sáng tối -- Sequential Blue): Biểu diễn độ lớn của giá trung bình (đậm = giá cao, nhạt = giá thấp).
    \end{itemize}
\end{itemize}

\vspace{0.2cm}
\textbf{Manipulate:}
\begin{itemize}[nosep, leftmargin=*, topsep=0.1cm]
    \item Navigate: Tích hợp tính năng Zoom/Pan giúp người dùng tự do phóng to, và di chuyển trên bản đồ để quan sát vi mô các phường nhỏ nằm sát nhau.
    \item Select: Cung cấp Tooltip để người dùng Hover chuột vào từng phân khu để xem chi tiết tên khu vực và con số AVG Price chính xác.
\end{itemize}

\vspace{0.2cm}
\textbf{Reduce:}
\begin{itemize}[nosep, leftmargin=*, topsep=0.1cm]
    \item Filter: Sử dụng bộ lọc Price Is Outlier (True/False/All) để kiểm soát việc hiển thị. Có thể chọn ``False'' để loại bỏ các căn siêu sang, giúp màu sắc bản đồ phản ánh đúng mức giá đại trà của thị trường.
\end{itemize} \\
\addlinespace[0.2cm]

\textbf{Scale} &
\begin{itemize}[nosep, leftmargin=*, topsep=0pt]
    \item Main key: Tọa độ không gian giới hạn trong khu vực NYC (hàng trăm phường/neighbourhoods).
    \item Color Range: Dải màu liên tục (Continuous) biểu diễn giá trị từ \$126.5 đến \$830.6.
\end{itemize} \\
\end{longtable}


\begin{figure}[H]
    \centering
    \safeincludegraphics[width=0.9\linewidth]{charts/domain_task1_chart1.png}
    \caption{Hình 1.1. Biểu đồ phân bố giá trung bình theo khu vực}
    \label{fig:domain-task1-chart1}
\end{figure}

\subsubsection*{B. Đánh giá biểu đồ}

\textbf{1. Nguyên lý biểu đạt (Expressiveness):}

\begin{itemize}
    \item Thuộc tính: Vị trí địa lý (Q), Khu vực (C), Giá trung bình (Q).
    \item Channel:
    \begin{itemize}
        \item PosX, PosY: định vị tọa độ địa lý $\rightarrow$ dùng cho thuộc tính không gian (Q).
        \item Color (Luminance -- độ đậm nhạt): thể hiện độ lớn của giá trị định lượng $\rightarrow$ dùng cho thuộc tính AVG Price (Q).
    \end{itemize}
    \item Đánh giá mức độ quan trọng:
    \begin{itemize}
        \item Ưu tiên 1 -- Vị trí không gian (PosX, PosY): Trong bản đồ, yếu tố cốt lõi nhất là định vị chính xác khu vực. Theo Mackinlay, kênh Vị trí (Position) là kênh mạnh nhất cho mọi loại dữ liệu. Việc gán tọa độ (Q) vào PosX/PosY đảm bảo tính chính xác tuyệt đối của ranh giới địa lý.
        \item Ưu tiên 2 -- Màu sắc (Color Luminance): Giá trị trung bình (Q) là thuộc tính quan trọng thứ hai. Do kênh Vị trí đã bị chiếm dụng cho tọa độ, thuộc tính này buộc phải sử dụng kênh Color Luminance (Độ sáng tối). Theo thang đo Mackinlay cho dữ liệu định lượng (Q), Color Luminance xếp hạng khá thấp, nhưng nó lại cực kỳ phù hợp cho mục đích Discover (nhìn ra bức tranh vĩ mô, xu hướng) thay vì so sánh con số tuyệt đối.
    \end{itemize}
    \item Kết luận: Dùng hoàn toàn chuẩn xác các kênh biểu đạt. Cấu trúc bản đồ tôn trọng tuyệt đối tính không gian thực tế của dữ liệu.
\end{itemize}

\textbf{2. Nguyên lý hiệu quả (Effectiveness):}

\begin{itemize}
    \item Độ chính xác (Accuracy):
    \begin{itemize}
        \item Kênh vị trí (PosX, PosY) có độ chính xác tuyệt đối theo dữ liệu địa lý.
        \item Tuy nhiên, kênh độ sáng màu (Color Luminance) biểu diễn dữ liệu định lượng (AVG Price) tuân theo định luật Stevens' Psychophysical law có sự sai lệch cảm nhận khá lớn, độ lỗi Log\_error rơi vào khoảng T9 $\leq 2.5$. Mắt người chỉ nhận biết được vùng nào đắt hơn (đậm hơn) dưới dạng thứ tự (Ordinal), chứ khó cảm nhận được lượng chênh lệch cụ thể là bao nhiêu USD. Thao tác Manipulate (Hover xem Tooltip) đã được thiết lập để bù đắp hoàn toàn nhược điểm này.
    \end{itemize}
    \item Khả năng phân biệt (Discriminability):
    \begin{itemize}
        \item Các khu vực (mark area) được chia theo đường ranh giới rõ ràng. Dải màu chuyển sắc (Sequential Blue) từ nhạt sang đậm có thể phân biệt được khoảng 5--7 cấp độ màu. Với các vùng liền kề có mức giá xấp xỉ nhau, màu sắc sẽ có xu hướng hòa vào nhau tạo thành các ``cụm khu vực đắt đỏ'' (ví dụ: cụm Manhattan bôi đậm), giúp mắt người nhìn ra xu hướng thay vì bị phân tâm bởi từng điểm nhỏ.
    \end{itemize}
    \item Khả năng tách biệt (Separability):
    \begin{itemize}
        \item Kênh vị trí (Pos) và kênh màu sắc (Color) tách biệt tốt, không bị can thiệp lẫn nhau. Vùng diện tích địa lý lớn hay nhỏ không làm sai lệch quá nhiều màu sắc tổng thể của vùng đó. Thao tác Reduce (Filter Price Is Outlier) giúp loại bỏ các giá trị nhiễu, giữ cho dải màu không bị bóp méo bởi một vài căn hộ giá hàng chục ngàn đô.
    \end{itemize}
\end{itemize}
